\documentclass{ximera}

\input{../../preamble.tex}

\title{the Unit Circle}

\begin{document}

\begin{abstract}
%Stuff can go here later if we want!
\end{abstract}
\maketitle




Sine and cosine are our descriptions of the coordinates on the unit circle. As a point moves counterclockwise on the unit circle, its coordinates change with the angle rotated.

We could think of two turns.

We could turn an angle. Then add on a second angle.

This forms three angles.

There is the first angle. There is the second angle. There is a third angle formed when the two separate angles are glued together.

Is there a connection between the sines and cosines of these three angles?





\subsection*{Learning Outcomes}

\begin{sectionOutcomes}
In this section, students will 

\begin{itemize}
\item explore the arithmetic of Complex Numbers.
\end{itemize}
\end{sectionOutcomes}















\begin{center}
\textbf{\textcolor{green!50!black}{ooooo-=-=-=-ooOoo-=-=-=-ooooo}} \\

more examples can be found by following this link\\ \link[More Examples of Trigonometric Forms]{https://ximera.osu.edu/csccmathematics/precalculus/precalculus/trigForms/examples/exampleList}

\end{center}






\end{document}
